\NormalFont\ShortTitle{Titus}

{\MT Scrisoarea lui Pavel către Titus


\par }\ChapOne{1}{\PP \VerseOne{1}Pavel, slujitor al lui Dumnezeu și apostol al lui Isus Hristos, după credința celor aleși de Dumnezeu și după cunoștința adevărului, care este potrivit cu evlavia,

\VS{2}în speranța vieții veșnice, pe care Dumnezeu, care nu poate minți, a făgăduit-o înainte de începutul veacurilor,

\VS{3}dar care, la vremea Sa, a descoperit cuvântul Său în mesajul care mi-a fost încredințat, potrivit poruncii lui Dumnezeu, Mântuitorul nostru,

\VS{4}lui Tit, adevăratul meu copil, potrivit unei credințe comune: Har, milă și pace de la Dumnezeu Tatăl și de la Domnul Isus Hristos, Mântuitorul nostru.


\par }{\PP \VS{5}De aceea v-am lăsat în Creta, ca să puneți în ordine ce lipsește și să rânduiți bătrâni în fiecare cetate, cum v-am poruncit eu,

\VS{6}dacă este cineva fără vină, soț al unei singure femei, cu copii credincioși, care nu sunt acuzați de desfrânare sau de neorânduială.

\VS{7}Pentru că supraveghetorul trebuie să fie fără vină, ca un administrator al lui Dumnezeu, să nu se mulțumească cu sine, să nu se mânie ușor, să nu fie dat la vin, să nu fie violent, să nu fie lacom la câștiguri necinstite;

\VS{8}ci să fie dat la ospitalitate, iubitor de bine, cumpătat, corect, cinstit, sfânt, stăpân pe sine,

\VS{9}să se țină de cuvântul credincios, care este conform învățăturii, ca să poată îndemna în doctrina sănătoasă și să convingă pe cei care îl contrazic.


\par }{\PP \VS{10}Căci sunt și mulți oameni nechibzuiți, vorbitori și înșelători, mai ales cei din circumcizie,

\VS{11}cărora trebuie să li se astupe gura, oameni care răstălmăcesc case întregi, învățând lucruri pe care nu trebuie să le facă, pentru un câștig necinstit.

\VS{12}Unul dintre ei, un profet de-al lor, a spus: “Creștinele sunt întotdeauna mincinoase, fiare rele și lăcomi leneși”.

\VS{13}Această mărturie este adevărată. De aceea, mustrați-i aspru, ca să fie sănătoși în credință,

\VS{14}fără să fie atenți la fabulațiile iudaice și la poruncile oamenilor care se îndepărtează de adevăr.

\VS{15}Pentru cei curați, toate lucrurile sunt pure, dar pentru cei necurați și necredincioși, nimic nu este curat, ci atât mintea lor, cât și conștiința lor sunt spurcate.

\VS{16}Ei mărturisesc că îl cunosc pe Dumnezeu, dar prin faptele lor îl neagă, fiind abominabili, neascultători și nepotriviți pentru orice lucrare bună.



\par }\Chap{2}{\PP \VerseOne{1}Ci spuneți ceea ce se potrivește cu învățătura sănătoasă,

\VS{2}ca bărbații mai în vârstă să fie cumpătați, cumpătați, cumpătați, cumpătați, sănătoși în credință, în dragoste și în perseverență,

\VS{3}și ca femeile mai în vârstă să fie, de asemenea, cuviincioase în purtare, să nu fie calomniatoare și să nu fie roabe de mult vin, învățătoare a ceea ce este bun,

\VS{4}pentru ca ele să învețe pe tinerele neveste să-și iubească soții, să-și iubească copiii,

\VS{5}să fie cumpătate, caste, lucrătoare în casă, amabile, supuse propriilor soți, pentru ca cuvântul lui Dumnezeu să nu fie hulit.


\par }{\PP \VS{6}De asemenea, îndemnați și pe cei mai tineri să fie cumpătați.

\VS{7}În toate lucrurile, arată-te pe tine însuți un exemplu de fapte bune. În învățătura voastră, arătați integritate, seriozitate, incoruptibilitate

\VS{8}și o vorbire sănătoasă, care nu poate fi condamnată, pentru ca cel care vi se opune să fie rușinat, neavând nimic rău de spus despre noi.


\par }{\PP \VS{9}Îndeamnă-i pe slujitori să se supună stăpânilor lor și să fie binevoitori în toate lucrurile, fără să se contrazică,

\VS{10}fără să fure, ci să dea dovadă de toată buna credincioșie, ca să împodobească în toate lucrurile învățătura lui Dumnezeu, Mântuitorul nostru.

\VS{11}Căci harul lui Dumnezeu s-a arătat, aducând mântuirea tuturor oamenilor,

\VS{12}instruindu-ne pentru ca, lepădându-ne de impietate și de poftele lumești, să trăim sobru, drept și evlavios în veacul acesta;

\VS{13}așteptând speranța binecuvântată și arătarea gloriei marelui nostru Dumnezeu și Mântuitor, Isus Cristos,

\VS{14}care s-a dat pe sine însuși pentru noi, ca să ne răscumpere de orice nelegiuire și să își purifice un popor ca să fie proprietatea sa, zelos pentru fapte bune.


\par }{\PP \VS{15}Spuneți aceste lucruri, îndemnați și mustrați cu toată autoritatea. Nimeni să nu vă disprețuiască.



\par }\Chap{3}{\PP \VerseOne{1}Amintește-le să se supună stăpânilor și autorităților, să fie ascultători, să fie gata pentru orice lucrare bună,

\VS{2}să nu vorbească de rău pe nimeni, să nu fie certăreți, să fie blânzi, cu toată smerenia față de toți oamenii.

\VS{3}Fiindcă și noi am fost odinioară nebuni, neascultători, înșelați, slujind diferitelor pofte și plăceri, trăind în răutate și invidie, urâcioși și urându-ne unii pe alții.

\VS{4}Dar, când s-a arătat bunătatea lui Dumnezeu, Mântuitorul nostru, și dragostea lui față de oameni,

\VS{5}nu prin fapte de dreptate pe care le-am făcut noi înșine, ci, potrivit milei sale, ne-a mântuit prin spălarea regenerării și înnoirea prin Duhul Sfânt,

\VS{6}pe care l-a revărsat din belșug asupra noastră prin Isus Hristos, Mântuitorul nostru,

\VS{7}pentru ca, fiind îndreptățiți prin harul său, să fim făcuți moștenitori, potrivit cu speranța vieții veșnice.

\VS{8}Această afirmație este credincioasă și, cu privire la aceste lucruri, doresc să insistați cu încredere, pentru ca cei care au crezut în Dumnezeu să fie atenți să păstreze faptele bune. Aceste lucruri sunt bune și folositoare pentru oameni;

\VS{9}dar evitați întrebările nebunești, genealogiile, certurile și disputele despre lege, pentru că sunt nefolositoare și zadarnice.

\VS{10}Evitați un om fățiș, după o primă și o a doua avertizare,

\VS{11}știind că unul ca acesta este pervertit și păcătos, fiind condamnat de sine.


\par }{\PP \VS{12}Când voi trimite la voi pe Artemas sau pe Tihicus, să vă sârguiți să veniți la mine la Nicopolis, căci am hotărât să iernez acolo.

\VS{13}Trimite repede pe Zenas, avocatul, și pe Apolo la drum, ca să nu le lipsească nimic.

\VS{14}De asemenea, poporul nostru să învețe să întrețină fapte bune pentru a satisface nevoile necesare, ca să nu rămână fără rodnicie.


\par }{\PP \VS{15}Toți cei ce sunt cu mine vă salută. Salutați-i pe cei care ne iubesc cu credință.

\par }{\PP Harul să fie cu voi toți. Amin.


\par }